% ==============================================================
%  Chapter 9: Do's and Don'ts
% ==============================================================
\chapter{Do's and Don'ts}
\label{ch:dosdontss}

This chapter collects the most important rules in one place.
They apply across all roles and are derived from what separates effective AI users
from ineffective ones --- and from what separates safe AI use from dangerous use.

\section{Universal Do's}
\label{sec:dos}

\begin{enumerate}

  \item \textbf{Give context before asking.}
    Tell AI who it is, what the goal is, what constraints apply, and what format you
    want. A well-framed prompt gets a useful answer; a vague prompt gets a vague
    answer.

  \item \textbf{Iterate, don't expect perfection first time.}
    The first output is a working draft. Follow up with \textit{``make it shorter''},
    \textit{``add an example''},\textit{``explain your reasoning''}.
    Iteration is the skill.

  \item \textbf{Verify before you use.}
    Read the output. Test the code. Check the facts against a primary source.
    AI produces plausible text, not verified truth.

  \item \textbf{Keep humans accountable for AI output.}
    Someone's name goes on everything AI helps produce. That person is responsible for
    its quality and accuracy.

  \item \textbf{Use AI to accelerate work you understand.}
    AI is a multiplier for existing skill, not a replacement for skill.
    The faster you can evaluate AI output, the more value you extract.

  \item \textbf{Ask for explanations, not just answers.}
    \textit{``Explain why''} and \textit{``what would break this?''} produce
    much deeper value than simply accepting the first answer.

  \item \textbf{Maintain a prompt library.}
    When a prompt produces a consistently useful result, save it.
    Share effective prompts with your team.

  \item \textbf{Update your mental model of AI regularly.}
    Models improve every few months. A limitation from six months ago may no longer
    apply. Stay current.

\end{enumerate}

\section{Universal Don'ts}
\label{sec:donts}

\begin{enumerate}

  \item \textbf{Don't share sensitive data with public AI tools.}
    This includes: customer PII, employee personal data, unpublished financial results,
    proprietary source code (unless using an enterprise API with data protection
    agreements), and active deal details.

  \item \textbf{Don't commit AI-generated code without understanding it.}
    If you cannot explain what every line does and why, you are not ready to commit it.
    This is a security risk, a quality risk, and a skills debt.

  \item \textbf{Don't use AI as a primary source for facts.}
    AI is a reasoning tool, not a research database. For factual claims, verify against
    primary sources (official docs, peer-reviewed papers, official statistics).

  \item \textbf{Don't skip review because AI said it looks fine.}
    AI reviewing its own output (or output in the same conversation) is not an
    independent check. Human review is still required.

  \item \textbf{Don't use AI for legal, medical, or financial decisions without expert review.}
    AI knows the concepts but cannot take professional responsibility. A lawyer,
    doctor, or accountant must own those decisions.

  \item \textbf{Don't use AI to generate credentials, secrets, or keys.}
    Any secret AI generates may exist in its context window or logs.
    Use a proper secrets manager.

  \item \textbf{Don't accept AI output that you cannot evaluate.}
    If you lack the domain knowledge to judge whether the output is correct, you need
    a human expert review --- not more AI confidence.

  \item \textbf{Don't measure AI adoption by volume.}
    Measuring prompts sent, documents generated, or code lines produced incentivises
    the wrong behaviour. Measure quality outcomes instead.

  \item \textbf{Don't let AI choose your values or make ethical judgements.}
    AI will produce an answer to any ethical question. That answer is pattern-matched
    from training data, not reasoned from principles. Ethical decisions belong to humans.

  \item \textbf{Don't expect AI to know your context without being told.}
    AI has no memory between sessions (unless explicitly provided), no knowledge of
    your organisation, and no access to your systems. Every conversation starts fresh.

\end{enumerate}

\section{Role-Specific Rules}
\label{sec:role-rules}

\begin{table}[htbp]
\centering
\caption{Critical rules by professional role}
\begin{tabularx}{\linewidth}{l X}
\toprule
\textbf{Role} & \textbf{Non-Negotiable Rule} \\
\midrule
Developer      & Never ship AI-generated code without manual review and test coverage \\[4pt]
DevOps         & Always run plan/dry-run before applying AI-generated IaC; check for permissive defaults \\[4pt]
Cloud Engineer & Verify AI-generated cost estimates against the current pricing calculator \\[4pt]
Data Scientist & Validate that AI-generated code does not introduce data leakage between train and test \\[4pt]
Teacher        & Personalise and verify all AI-generated feedback before sending to students \\[4pt]
Manager        & Never share customer PII, employee data, or confidential strategy in public AI tools \\[4pt]
All roles      & Take full accountability for any AI output that carries your name \\
\bottomrule
\end{tabularx}
\label{tab:role-rules}
\end{table}

\section{Security and Privacy Non-Negotiables}
\label{sec:security-privacy}

These apply to everyone, regardless of role:

\begin{itemize}
  \item Use enterprise AI subscriptions (not free tiers) for any work-related data.
        Enterprise tiers typically include data processing agreements and opt-out of
        training use.
  \item Never paste passwords, API keys, private SSH keys, or OAuth tokens into any
        AI tool.
  \item Treat AI prompts as potentially logged. Write prompts as if they could be read
        by your company's security team.
  \item Report AI-related data incidents through your normal incident reporting channel.
        Do not try to resolve them silently.
\end{itemize}
